\documentclass[aspectratio=1610,12pt,xcolor=dvipsnames]{beamer}
\mode<presentation>

% Respect chosen fonts (serif) and metrics
\usefonttheme{professionalfonts}
\renewcommand{\familydefault}{\rmdefault}

% Figures
\usepackage{graphicx}
\usepackage{float}
\usepackage{tikz}
\usetikzlibrary{arrows.meta, positioning, fit}

%% table
\usepackage{xcolor}
\usepackage{fancyhdr}
\usepackage{booktabs}
\usepackage[table]{xcolor}
\usepackage{tabularx,array}

% subsection page template
\setbeamertemplate{subsection page}
{
  \begin{centering}
    \vfill
    {\usebeamerfont{section title}\usebeamercolor[fg]{section title}\Large\insertsubsectionhead\par}
    \vfill
  \end{centering}
}
% section page template
\setbeamertemplate{section page}
{
  \begin{centering}
    \vfill
    {\usebeamerfont{section title}\usebeamercolor[fg]{section title}\Large\insertsectionhead\par}
    \vfill
  \end{centering}
}

%% footnote
\setbeamertemplate{footnote}{%
  \parindent 0em\noindent%
  \raggedright\insertfootnotemark\insertfootnotetext\par%
}

%% theme
\usetheme{default}
\useoutertheme{miniframes}
\definecolor{nagivation}{rgb}{0,0,0.35} % (typo kept if intentional)
\definecolor{main}{rgb}{0,0,0.5}
\setbeamercolor{structure}{bg=white,fg=nagivation}
\setbeamercolor{frametitle}{fg=nagivation}
\setbeamercolor{section in head/foot}{fg=white,bg=nagivation}
\setbeamertemplate{footline}[page number]
\setbeamertemplate{items}[circle]
\setbeamerfont{footnote}{size=\footnotesize}
\setbeamerfont{caption}{size=\footnotesize}
\setbeamertemplate{subsection in head/foot}{}%
\setbeamertemplate{subsection in head/foot shaded}{}%

%% font
\usepackage{dsfont}
\usepackage[T1]{fontenc}
\usepackage{lmodern}
\usepackage{amsmath}

\setbeamerfont{section title}{series=\bfseries}
\setbeamerfont{subsection title}{series=\bfseries}
\setbeamerfont{block title}{series=\bfseries}

% Math & bib
\usepackage{amsmath,amsfonts,amssymb,bm}
\DeclareMathOperator*{\argmin}{arg\,min}
\newcommand{\indep}{\perp\!\!\!\, \perp}
\usepackage{natbib}
\bibliographystyle{asr}
\setcitestyle{aysep={}}
\usepackage[english]{babel}

\makeatletter
\DeclareRobustCommand\citep
{\begingroup\scriptsize\color{gray}\NAT@swatrue\let\NAT@ctype\z@\NAT@partrue
    \@ifstar{\NAT@fulltrue\NAT@citetp}{\NAT@fullfalse\NAT@citetp}}
\makeatother

% Roman numerals macro
\newcommand{\rom}[1]{\uppercase\expandafter{\romannumeral #1\relax}}

%% Title
\title[CML]{SOC 690S: Machine Learning in Causal Inference\\[1.5pt]}
\subtitle{\large Week 14: Feature Engineering for Causal Inference\footnote{\scriptsize I heavily draw upon the materials from \href{https://causalml-book.org/}{Chapter 10} of the textbook of Causal Machine Learning by Chernozhukov et al. (2026)} \\[-10pt]}

%% author
\author[Jiang]
{\large Wenhao Jiang\vspace{-2em}}

%% affiliation
\institute[Duke]{}
\titlegraphic{\includegraphics[height=1.4cm]{Misc/duke_logo.png}}

\date[Duke]
{\large Department of Sociology, Fall 2025}

\begin{document}

%%%%%%%%%%%%%%%%%%%%%%%%%%%%%%%%%%%%
%%%%%%%%Begin Main Content%%%%%%%%%%
%%%%%%%%%%%%%%%%%%%%%%%%%%%%%%%%%%%%

%% Title page %%
\begin{frame}
    \titlepage 
\end{frame}

\section{Introduction and Motivation}

\begin{frame}{Introduction and Motivation}

\begin{itemize}
    \item Social scientists increasingly leverage unstructured, multimodal data to study human behavior, social interactions, and cultural and institutional processes \citep{lazer2009computational}
    
    \item These data are often represented as vector \textit{embeddings} for machines to understand and analyze \pause
    
    \item Sociologists have primarily used embeddings in two ways:
    \begin{itemize}
        \item As \textit{objects of study} that capture cultural meaning, symbolic associations, and social representations \citep{kozlowski2019geometry}
        \item As predictive features that improve downstream tasks, often by fine-tuning pre-trained models \citep{bonikowski2022politics}
    \end{itemize} \pause
    \item Despite their substantive richness, embeddings are rarely incorporated as \textit{features} into standard causal inference, even though they may encode important confounding information that should be adjusted for in regression analyses
\end{itemize}

\end{frame}

\begin{frame}{Introduction and Motivation}

\begin{itemize}
    \item Causal inference researchers are interested in estimating the treatment effect of $D$ on the outcome $Y$
    \[ Y = \beta_1 D + \beta_2' W + \epsilon \]
    \item In observational settings, we rely on the \textit{strong ignorability} assumption, which requires that all confounders are observed and properly modeled by $\beta_2' W$ \pause
    \item This is a strong assumption because there are numerous potential confounders, some of which are difficult to measure using, for example, survey data

\end{itemize}
\end{frame}

\begin{frame}{Introduction and Motivation}

\begin{itemize}
    \item Consider the \textit{neighborhood effects} literature, where neighborhood context is argued to have a causal effect on individual outcomes
    \item For example, in studying the effect of gunshot presence in a neighborhood on children's academic performance, one may need to control for overall neighborhood ``quality''
    \item Increasingly, researchers use data such as images and text to measure neighborhood ``quality,'' which is often not captured in survey data
\end{itemize}
\end{frame}

\begin{frame}{Introduction and Motivation}

\begin{itemize}
    \item Consider estimating the \textit{price elasticity of demand}: the causal effect of price on consumer demand
    \item Comparing products that differ substantially in quality, design, or target market can confound this relationship
    \item Credible estimation requires controlling for product quality and other characteristics, but these attributes are often difficult to measure using conventional structured data \citep{bajari2025hedonic} \pause
    \begin{itemize}
        \item However, they may be encoded in product images or text as embeddings
    \end{itemize} 
\end{itemize}
    
\end{frame}

\begin{frame}{Overview of the Talk}
\begin{itemize}
    \item Introduce several approaches for transforming complex data—particularly text, which encodes rich information—into low-dimensional representations   
    \item Review Double Machine Learning (DML), which enables flexible adjustment for high-dimensional feature representations in regression\footnote{\scriptsize ``Low-dimensional'' is relative to the original text data, which may contain millions of unique words and their relationships; ``high-dimensional'' is relative to the sample size ($p/n \not\rightarrow 0$)}
    \begin{itemize}
        \item Brief review of OLS, the Frisch--Waugh--Lovell (FWL) theorem, and their asymptotic properties
        \item Introduction to \textit{Neyman Orthogonality}
    \end{itemize}
    \item Illustrate the methods  using the example of estimating the price elasticity of demand from Amazon's data
\end{itemize}
\end{frame}

\section{Double Machine Learning (DML)}

\begin{frame}
  \sectionpage
\end{frame}

\begin{frame}{OLS and the Frisch--Waugh--Lovell (FWL) Theorem}

\begin{itemize}
    \item Consider the linear regression model under low-dimensional regime
    \begin{align*}
        Y = \beta_1 D + \beta_2'W + \epsilon
    \end{align*}

    \item FWL Theorem states that the coefficient $\beta_1$ can be obtained by regressing the residualized outcome on the residualized treatment\footnotetext{\scriptsize For a proof, see p.~32 week 1 slides: \url{https://github.com/wenhaojiangsoc/CML/tree/main}}
    \begin{align*}
        \beta_1
        =
        \arg\min_{b}
        E\!\left[(\tilde Y-b\tilde D)^2\right]
    \end{align*} \\[-2em]
    \begin{align*}
        \tilde D &= D-\gamma_D'W;
        &\hspace{-3em}
        \gamma_D
        &=\arg\min_{\gamma}
        E\!\left[(D-\gamma'W)^2\right] \\
        \tilde Y &= Y-\gamma_Y'W;
        &\hspace{-3em}
        \gamma_Y
        &=\arg\min_{\gamma}
        E\!\left[(Y-\gamma'W)^2\right]
    \end{align*}
\end{itemize}
\end{frame}

\begin{frame}{OLS and the Frisch--Waugh--Lovell (FWL) Theorem}

\begin{itemize}
    \item Under regularity conditions and low-dimensional data $p/n \rightarrow 0$
    \begin{align*}
    \sqrt{n}\left(\hat{\beta}_1-\beta_1\right)
    \xrightarrow{d}
    \mathcal{N}\!\left(
    0,\,
    E\!\left[\tilde D\tilde D'\right]^{-1}
    E\!\left[\epsilon^2\tilde D\tilde D'\right]
    E\!\left[\tilde D\tilde D'\right]^{-1}
    \right)
\end{align*}
    \item $\hat \beta_1$ is $\sqrt n$-consistent
    \item The consistent ``sandwich'' estimator (Eicker-Huber-White) of a \textit{sample} is then given by
\begin{align*}
    \hat{V}(\hat{\beta}_1) = \left(\tilde D \tilde D'\right)^{-1}\left( \sum_i^{n} \tilde D \tilde D'\hat{\epsilon}^{2} \right) \left(\tilde D \tilde D'\right)^{-1}
\end{align*}
by plugging in sample $\hat \epsilon^2$ to estimate $\epsilon^2$
\end{itemize}
\end{frame}

\begin{frame}{Double Machine Learning (DML)}

\begin{itemize}
    \item Under high-dimensional regime $p/n \not\rightarrow 0$, the nuisance regressions can no longer be estimated precisely enough for conventional OLS inference; $\sqrt n$ asymptotic normality generally fails

    \item Double Machine Learning (DML) yields $\sqrt n$-consistent estimates and valid standard errors under \textit{Neyman Orthogonality} and cross-fitting
\end{itemize}
\end{frame}

\begin{frame}{Double Machine Learning (DML)}

\begin{itemize}
    \item Consider a partially linear model under a high-dimensional regime
    \begin{align*}
        Y &= \beta_1 D + \eta(W) + \epsilon \\
        \beta_1 &=
        \arg\min_{b}
        E\!\left[(\tilde Y-b\tilde D)^2\right] \\
        \tilde Y &= Y - E\!\left[ Y \mid W \right] \\
        \tilde D &= D - E\!\left[ D \mid W \right]
    \end{align*}

    \item The nuisance functions $g(W)=E\!\left[Y\mid W\right]$ and $m(W)=E\!\left[D\mid W\right]$ can be estimated flexibly using nonlinear models.

    \item Machine learning provides a natural class of estimators for these nuisance functions when $W$ is high-dimensional.
\end{itemize}

\end{frame}

\begin{frame}{Cross-Fitting in Double Machine Learning (DML)}

\begin{itemize}
    \item Nuisance functions and residualization are operationalized through \textit{cross-fitting}
    \begin{figure}
        \centering
        \includegraphics[width=0.67\linewidth]{Feature Engineering/Figures/cross-fitting.pdf}
    \end{figure}
    \item The nuisance functions $g(W) = E\!\left[Y \mid W\right]$ and $m(W) = E\!\left[D \mid W\right]$ are estimated in the \textcolor{blue!45}{blue fold}, and $\tilde Y$ and $\tilde D$ are estimated in the \textcolor{red!45}{red fold}
\end{itemize}
    
\end{frame}

\begin{frame}{Neyman Orthogonality in Double Machine Learning (DML)}

\begin{itemize}

\item The key property underlying Double Machine Learning is \textit{Neyman Orthogonality} through \textit{cross-fitting}

\item Consider the moment function
\begin{align*}
\psi \left(Y,D,W;\beta,g,m \right) &=
(\tilde Y-\beta\tilde D)\tilde D \\
\tilde Y = Y-g(W), &\:\:
\tilde D = D-m(W) \\[-2.6em]
\end{align*} \pause
\item The moment is Neyman orthogonal if its (Gateaux) derivative with respect to the nuisance functions vanishes at the truth
\[
\left.
\frac{\partial E\!\left[\psi(Y,D,W;\beta,g,m)\right]}{\partial(g,m)}
\right|_{(g,m)=(g_0,m_0)}
=0
\]

\item Intuitively, the moment function is locally \emph{flat} in the nuisance direction

\item Consequently, small errors in estimating $g$ and $m$ contribute only \emph{second-order} bias, allowing valid $\sqrt n$-inference for $\beta_1$
\end{itemize}
\end{frame}

\begin{frame}{Double Machine Learning: Neyman Orthogonality}

\begin{figure}
    \centering
    \hspace*{-1.6cm} % shift left
    \vspace*{2em} % shift up
    \includegraphics[width=1.15\linewidth]{Causal Inference Potential Outcome/Figures/orthogonal vs. non.png}
\end{figure}
\end{frame}

\begin{frame}{Double Machine Learning with Embeddings}

\begin{itemize}
    \item We estimate the effect of product price on consumer demand while controlling for high-dimensional covariates $W$, such as text or image embeddings that encode rich information about product characteristics
    
    \item Rather than assuming a fully linear model, we consider the partially linear model
    \begin{align*}
        Y = \beta_1 D + \eta(W) + \epsilon
    \end{align*}
    
    \item The target parameter $\beta_1$ is estimated by
    \[
        \hat \beta_1
        =
        \arg\min_{b}
        E\!\left[(\tilde Y-b\tilde D)^2\right]
    \]
    where $\tilde Y$ and $\tilde D$ are residualized using flexible machine learning methods with cross-fitting
\end{itemize}
\end{frame}

\section{Feature Extraction}

\begin{frame}
  \sectionpage
\end{frame}

\begin{frame}{Feature Extraction: Principal Components}

\begin{itemize}
    \item We are interested in transforming complex objects such as text into a collection of relatively low-dimensional numerical features
    \item Principal components are an early classic example of embeddings
    \item One way to frame principal components is that principal components find orthonormal linear combinations that are ``best'' at reproducing the underlying data \pause
    \item A small number of principal components should capture most of the variability in the original variables and provide a useful low-dimensional summary of the original data
\end{itemize}
\end{frame}

\begin{frame}{Feature Extraction: Principal Components}

\begin{itemize}
    \item Let $(W_1, ..., W_n)$ be a sample of $n$ observations of a high-dimensional centered random vector $W$ in $\mathds{R}^d$, and let $\Sigma_n = \mathbb{E}\left[WW'\right] \in \mathds{R}^{d \times d}$ denote the empirical covariance matrix \pause
    \item The eigenvector $c_k$ in $\mathds{R}^d$ and eigenvalue $\lambda_k$ is found by 
    \begin{align*}
        \Sigma_n c_k = \lambda_k c_k
    \end{align*}
    \item We find the $K$ ($K \ll d$) largest eigenvalues in an descending order and corresponding eigenvectors $C_K = [c_1,...,c_K]$
    \item We find $W_i$'s projection on each eigenvector (principal direction)
    \begin{align*}
        X_{ki} = c_k' W_i, \quad k = 1, ..., K
    \end{align*}
    \item Principal components $X_{i}^{K} \in \mathds{R}^{K}$
    \begin{align*}
        X_{i}^{K} = (X_{1i},...,X_{Ki})'
    \end{align*}
\end{itemize}
\end{frame}

\begin{frame}{Feature Extraction: Principal Components}

\begin{itemize}
    \item Principal components $X_{i}^{K} \in \mathbb{R}^{K}$
    \begin{align*}
        X_{i}^{K} = (X_{1i},...,X_{Ki})'
    \end{align*}
    \item $X_i^K$ represents the compressed feature vector of the original high-dimensional vector $W$ 
    \begin{figure}
    \centering
    \includegraphics[width=0.95\linewidth]{Feature Engineering/Figures/PCA_illustration.pdf}
    \end{figure}
\end{itemize}
\end{frame}

\begin{frame}{Feature Extraction: Principal Components}

\begin{itemize}
    \item Principal components $X_{i}^{K} \in \mathbb{R}^{K}$
    \begin{align*}
        X_{i}^{K} = (X_{1i},...,X_{Ki})'
    \end{align*}
    \item $X_i^K$ represents the compressed feature vector of the original high-dimensional vector $W$ 
    \item One common use of Principal Components is to compress image
    \item We approximate the original vector $W_i$ by projecting $X_i^K$ back to the $d$-dimensional space
    \[
    \hat{W}_i = C_K X_i^K
    \]
    \item The choice of $K$ determines how much of the original information (variance) in $W_i$ is retained
\end{itemize}


\end{frame}

\begin{frame}{Principal Components as Image Compression}

\begin{figure}
\centering

\begin{minipage}{0.34\textwidth}
    \centering
    \includegraphics[width=\linewidth]{Feature Engineering/Figures/pca_fig1_matrix.png}
\end{minipage}
\hfill
\begin{minipage}{0.32\textwidth}
    \centering
    \includegraphics[width=\linewidth]{Feature Engineering/Figures/pca_fig2_10pc.png}
\end{minipage}
\hfill
\begin{minipage}{0.32\textwidth}
    \centering
    \includegraphics[width=\linewidth]{Feature Engineering/Figures/pca_fig3_50pc.png}
\end{minipage}
\end{figure}
\end{frame}

\begin{frame}{Principal Components in Other Applications}

\begin{itemize}
    \item Principal Component Analysis (PCA), or equivalently Singular Value Decomposition (SVD), underlies many important applications
    \item Latent Semantic Analysis (LSA) creates low-dimensional embeddings for words and documents from a word--document matrix
    \item Recommender systems create low-dimensional embeddings for users and items from a user--item interaction matrix
\end{itemize}
    
\end{frame}

\begin{frame}{Principal Components Projection as Autoencoder}

\begin{itemize}
    \item The approximation to $W_i$ from $W$ can be seen as a predictive exercise
    \item It underlies a shallow one-layer \textit{linear} neural network
    \begin{align*}
    \underset{d \times 1}{W_i}
    \;\mapsto\;
    C_K' W_i := \underset{K \times 1}{E_i}
    \;\mapsto\;
    A' E_i := \underset{d \times 1}{\hat{W}_i}
    \end{align*}
    \item The first step is said to be ``encoding'' the information in the input, and the second step is said to be ``decoding'' in the sense of returning the encoded information to the original space
    \item Principal components projections are embeddings generated by a linear ``encoder--decoder'' network (an \textit{autoencoder})
    \item For principal components projection, $A = C_K'$
\end{itemize}
\end{frame}

\begin{frame}{Principal Components as Autoencoder}

\begin{minipage}{0.5\textwidth}
    \centering
    \includegraphics[width=\linewidth]{Feature Engineering/Figures/autoencoder1.pdf}
\end{minipage}
\hfill
\begin{minipage}{0.45\textwidth}
    \begin{itemize}
        \item The weights connecting the input to each hidden neuron correspond to a principal direction (eigenvector)
        \item The encoding layer outputs are the principal component
        \item The decoder reconstructs the input using these components
    \end{itemize}
\end{minipage}
\end{frame}

\begin{frame}{Principal Components as Autoencoder}

\begin{minipage}{0.5\textwidth}
    \centering
    \includegraphics[width=\linewidth]{Feature Engineering/Figures/autoencoder1_colored.pdf}
\end{minipage} \pause
\hfill
\begin{minipage}{0.45\textwidth}
    \begin{itemize}
        \item The weights connecting each input neuron to hidden neurons correspond to each feature's weight or coefficient on each latent dimension
    \end{itemize}
\end{minipage}
\end{frame}

\begin{frame}{From Principal Components To Autoencoder}
    
\begin{minipage}{0.7\textwidth}
    \centering
    \includegraphics[width=\linewidth]{Feature Engineering/Figures/autoencoder2.pdf}
\end{minipage}
\hfill
\begin{minipage}{0.28\textwidth}
    \begin{itemize}
        \item The approach can be generalized to nonlinear encoders and decoders with multiple layers with \textit{nonlinear} activation functions, forming deep autoencoders
    \end{itemize}
\end{minipage}
    \[
    W_i \;\overset{g_1}{\mapsto}\; E_i^{1} \;\cdots\; \overset{g_k}{\mapsto}\; E_i^{k}
    \;\overset{g_{k+1}}{\mapsto}\; D_i^{k+1} \;\cdots\; \overset{g_m}{\mapsto}\; D_i^{m}
    \;:=\; \hat{W}_i
    \]   
\end{frame}

\begin{frame}{From Principal Components To Autoencoder}
    
\begin{minipage}{0.7\textwidth}
    \centering
    \includegraphics[width=\linewidth]{Feature Engineering/Figures/autoencoder2.pdf}
\end{minipage}
\hfill
\begin{minipage}{0.28\textwidth}
    \begin{itemize}
        \item Deep autoencoders preserve more information in the original data
        \item The trained encoder (the weights) is used to generate embeddings for downstream tasks
    \end{itemize}
\end{minipage}
    \[
    W_i \;\overset{g_1}{\mapsto}\; E_i^{1} \;\cdots\; \overset{g_k}{\mapsto}\; E_i^{k}
    \;\overset{g_{k+1}}{\mapsto}\; D_i^{k+1} \;\cdots\; \overset{g_m}{\mapsto}\; D_i^{m}
    \;:=\; \hat{W}_i
    \]   
\end{frame}

\begin{frame}{From Autoencoders to General Embeddings}

\begin{itemize}
    \item The encoder--decoder paradigm underlies many modern representation learning methods, including deep autoencoders and encoder--decoder Transformers such as T5 and BART
    
    \item In a classical autoencoder, the objective is to minimize the expected reconstruction loss
    \[
        \mathbb{E}\!\left[\mathrm{loss}\!\left(W,\; d(e(W))\right)\right]
    \]
    where $e(\cdot)$ and $d(\cdot)$ denote the encoder and decoder

    \item More generally, the prediction target need not be the original input
    \[
        \mathbb{E}\!\left[\mathrm{loss}\!\left(A,\; f(e(W))\right)\right]
    \]
    where $A$ is an arbitrary learning target
\end{itemize}
\end{frame}

\begin{frame}{From Autoencoders to General Embeddings}

\begin{itemize}
    \item More generally, the prediction target need not be the original input
        \[
            \mathbb{E}\!\left[\mathrm{loss}\!\left(A,\; f(e(W))\right)\right]
        \]
        where $A$ is an arbitrary learning target
    \item For example, BERT predicts masked words in a sentence, while T5 and BART reconstruct corrupted text from its encoded representation

    \item The resulting embeddings often capture more useful semantic information than simple reconstruction

    \item In our price elasticity example from Amazon products we will see later, we may prefer embeddings that encode product type and functionality rather than faithfully reconstruct the product description
\end{itemize}
    
\end{frame}

\section{Word2Vec}

\begin{frame}
  \sectionpage
\end{frame}

\begin{frame}{Word2Vec Embeddings}

\begin{itemize}
    \item Different learning objectives $A$ produce different types of embeddings
    \item Word2Vec learns embeddings by predicting a word from its surrounding context (or the other way around)
    \item This idea is motivated by the \textit{distributional hypothesis}: words are understood by their semantic context
\end{itemize}

\vspace{0.0em}

\begin{minipage}{0.55\textwidth}
\begin{itemize}
    \item A bottle of \textcolor{blue}{tesgüino} is on the table
    \item Everyone likes \textcolor{blue}{tesgüino}
    \item \textcolor{blue}{Tesgüino} makes you drunk
    \item We make \textcolor{blue}{tesgüino} out of corn
\end{itemize}
\end{minipage}
\hspace{2em}
\begin{minipage}{0.23\textwidth}
\centering
\pause
\includegraphics[width=\linewidth]{Feature Engineering/Figures/tesguino.jpg}
\end{minipage}
\end{frame}

\begin{frame}{Word2Vec Embeddings}

\begin{itemize}
    \item Two training objectives are commonly used in Word2Vec

    \begin{itemize}
        \item \textit{Continuous Bag-of-Words (CBOW):} predict the center word from its surrounding context
        \item \textit{Skip-gram:} predict the surrounding context words from the center word
    \end{itemize}
\end{itemize}

\end{frame}

\begin{frame}{Word2Vec: Skip-gram with Negative Sampling}

\begin{minipage}{0.62\textwidth}
\centering
\includegraphics[width=\linewidth]{Feature Engineering/Figures/skipgram.pdf}
\end{minipage}%%
\begin{minipage}{0.37\textwidth}
\begin{itemize}
    \item Analogous to Principal Components Projection, each embedding ``weight'' (from input to embedding layer) learns the coefficient of the word to each latent semantic coordinate from the word co-occurrence structure
\end{itemize}
\end{minipage}
\end{frame}

\begin{frame}{Word2Vec Embeddings}

\begin{minipage}{0.4\textwidth}
\centering
\includegraphics[width=\linewidth]{Feature Engineering/Figures/word_analogy.jpg}
\end{minipage}%%
\begin{minipage}{0.6\textwidth}
\begin{itemize}
    \item Word2Vec learns a low-dimensional vector representation for each word
    \item Words appearing in similar contexts tend to have similar embeddings
    \item The learned embeddings capture semantic relationships, such as word analogies \pause
    \item However, the embeddings are context-independent: ``bank'' in ``river bank'' and ``financial bank'' receive the same representation
    \item Word order and syntax are not explicitly modeled; sentence embeddings formed by averaging word embeddings are often imprecise
\end{itemize}
\end{minipage}
\end{frame}

\section{ELMo}

\begin{frame}
  \sectionpage
\end{frame}

\begin{frame}{Sequence Models}

\begin{itemize}
    \item A major advance beyond Word2Vec was to represent text as a \textit{sequence}, allowing embeddings to depend on word order and surrounding context
    
    \item For example, ``CASER belongs to NYU Shanghai'' conveys a different meaning than ``NYU Shanghai belongs to CASER'' \pause
    
    \item One influential approach is \textit{Embeddings from Language Models} (ELMo), which learns context-dependent word embeddings
    
    \item More generally, language models are inspired by the Shannon game: predict a missing or future word given the surrounding words in a sentence
\end{itemize}

\end{frame}

\begin{frame}{Sequence Model: ELMo}

\begin{itemize}
    \item Let $T_k$ denote the $k$-th token in a sentence of length $n$, and let $t$ denote a candidate word in the vocabulary
    \item The forward language model predicts the next word using the preceding words
    \[
        p_k^{f}(t)
        =
        P\!\left(
        T_{k+1}=t
        \mid
        T_1,\ldots,T_k;\theta
        \right)
    \]
    \item The backward language model predicts the previous word using the following words:
    \[
        p_k^{b}(t)
        =
        P\!\left(
        T_{k-1}=t
        \mid
        T_k,\ldots,T_n;\theta
        \right)
    \]
    \item Here $\theta$ denotes all model parameters
\end{itemize}
\end{frame}

\begin{frame}{Sequence Model: ELMo}

\begin{itemize}
    \item Many social and economic processes are inherently sequential; for example, life-course trajectories in sociology and financial time series depend on past histories

    \item Likewise, language exhibits sequential dependence: the probability of the next word depends on the preceding words

    \item Recurrent Neural Networks (RNNs) model such sequential dependencies by maintaining a hidden state that evolves over the sequence

    \item To mitigate vanishing and exploding gradients over long sequences, ELMo employs Long Short-Term Memory (LSTM) networks, whose gated architecture enables information to persist across many time steps
\end{itemize}

\end{frame}

\begin{frame}{Sequence Model: ELMo}

\begin{minipage}{0.65\textwidth}
\centering
\includegraphics[width=\linewidth]{Feature Engineering/Figures/ELMo.pdf}
\end{minipage}%%
\begin{minipage}{0.34\textwidth}
\begin{itemize}
    \item Input $x_i$ consists of context-free word embeddings
    \item Two LSTM networks are used, forward and backward
    \item After pre-training, each word is associated with $2L$ contextual hidden states, where $L$ is the number of stacked LSTM layers in each direction
\end{itemize}
\end{minipage}
\end{frame}

\begin{frame}{Sequence Model: ELMo: ELMo and Downstream Task}

\begin{minipage}{0.65\textwidth}
\[t_k \mapsto w_k := \sum_{i=1}^{L} (\gamma_i w^f_{ki} + \bar{\gamma}_i w^b_{ki})\]
\centering
\includegraphics[width=\linewidth]{Feature Engineering/Figures/ELMo.pdf}
\end{minipage}%%
\begin{minipage}{0.34\textwidth}
\begin{itemize}
    \item $2L$ hidden states are combined to produce the final context-dependent ELMo embedding
    \item The weights $\gamma$ and $\bar \gamma$ are learned during the downstream task (e.g., classification) to combine hidden states into context-dependent word representations
\end{itemize}
\end{minipage}
\end{frame}

\section{BERT}

\begin{frame}
  \sectionpage
\end{frame}

\begin{frame}{Transformer Architecture}

\begin{itemize}
    \item A recent advance in language modeling is the transformer architecture
    
    \item Rather than processing text sequentially, transformers use \textit{self-attention} to model the relationships among all words in a sentence simultaneously
    
    \item Compared with LSTMs, self-attention more effectively captures long-range dependencies by allowing each word to attend directly to the most relevant words, regardless of their distance in the sequence \pause
    \item For example,\\[0.5em]
    \begin{quote} 
        \small{The Summer Institute in Computational Social Science (SICSS) at NYU Shanghai was hosted by CASER. \textbf{It} attracts smart audience from around the world.}
    \end{quote} 
    
    \item Self-attention helps determine that ``\textit{it}'' refers to ``SICSS'', even though the relevant words are far apart
\end{itemize}

\end{frame}

\begin{frame}{Transformer Architecture}

\begin{itemize}
    \item A prominent transformer-based language model is \textit{Bidirectional Encoder Representations from Transformers} (BERT)

    \item The original BERT model is pre-trained using two self-supervised objectives

    \begin{itemize}
        \item \textit{Masked Language Modeling (MLM):} Randomly mask a subset of words in a sentence and predict the masked words from their surrounding context

        \item \textit{Next Sentence Prediction (NSP):} Given a pair of sentences, predict whether the second sentence immediately follows the first in the original text
    \end{itemize}
\end{itemize}
\end{frame}

\begin{frame}{Transformer Architecture: BERT\footnote{\scriptsize Check out this site to calculate attention by hand \url{https://www.byhand.ai/p/11-can-you-calculate-self-attention}}}

\begin{itemize}

\item For a sentence of $n$ tokens, each token is represented by the sum of its token embedding (e.g., Word2Vec) and positional embedding. Collecting all token representations gives
\[ X \in \mathds{R}^{n \times d}
\]
\item Query, Key, and Value matrices are computed by
\begin{align*}
    Q &= XW^Q \in \mathds{R}^{n \times d_k} \\
    K &= XW^K \in \mathds{R}^{n \times d_k} \\
    V &= XW^V \in \mathds{R}^{n \times d_v} \\
    W^Q,\;W^K &\in \mathds{R}^{d \times d_k} \\
    W^V &\in \mathds{R}^{d \times d_v}
\end{align*}
\end{itemize}
\end{frame}

\begin{frame}{Transformer Architecture: BERT}

\begin{itemize}
    \item Self-attention is computed as
    \[
    \operatorname{Attention}(Q,K,V)
    =
    \operatorname{softmax}
    \!\left(
    \frac{QK'}{\sqrt{d_k}}
    \right)V
    \]
    \item Since $QK' \in \mathds{R}^{n\times n}$, each token computes an attention score for every token in the sequence, including itself

    \item The resulting $n\times n$ attention matrix determines how information is aggregated across all tokens to produce contextual embeddings
    \begin{itemize}
        \item Each entry asks how much should token $i$ attend to token $j$
    \end{itemize}
    \item The rows of $\operatorname{Attention}(Q,K,V) \in \mathds{R}^{n \times d_v}$ are contextual representations of the tokens, which are further refined by the remainder of the Transformer block
\end{itemize}
\end{frame}

\begin{frame}{Transformer Architecture: BERT}

\begin{itemize}
    \item During pre-training, the projection matrices $W^Q$, $W^K$, and $W^V$ (together with all other model parameters) are updated by gradient descent, and $W^Q$, $W^K$, and $W^V$ are global model parameters

    \item For example, consider the masked sentence
    \[
    \texttt{[CLS] The SICSS at NYU Shanghai was hosted by [MASK] . [SEP]}
    \]

    \item The model computes a contextual embedding for every token, uses the contextual embedding of \texttt{[MASK]} to predict the missing word, and compares the prediction with the true word (e.g., \texttt{CASER})

    \item The prediction error is backpropagated to update $W^Q$, $W^K$, $W^V$, the feed-forward layers, and the token embeddings
    
    \item After training, the encoder (e.g., projection matrices, token embedding, residual connections) is retained to generate contextual embeddings
\end{itemize}
\end{frame}

\begin{frame}{Transformer Architecture: BERT}

\begin{minipage}{0.6\textwidth}
\centering
\includegraphics[width=\linewidth]{Feature Engineering/Figures/BERT.pdf}
\end{minipage}%%
\begin{minipage}{0.4\textwidth}
\begin{itemize}

    \item The last Transformer layer produces the final hidden states $h_i^{(L)}$, which serve as the contextual embedding for each token

    \item A sentence embedding can then be obtained by aggregating the token embeddings (e.g., averaging or pooling) or, in BERT, by using the contextual embedding of the special \texttt{[CLS]} token

\end{itemize}
\end{minipage}
\end{frame}

\begin{frame}{Fine-Tuning BERT for Downstream Tasks}
\begin{itemize}     
    \item There are two common approaches to adapting a pre-trained BERT model for downstream tasks    
    \item \textit{Linear probing}: Freeze the pre-trained BERT encoder and train only a linear prediction head on top of the final hidden states         
    \item \textit{Fine-tuning}: Update both the BERT encoder and the task-specific prediction head jointly using the downstream training data
    \end{itemize} 
\end{frame}

\begin{frame}{Fine-Tuning BERT for Nuisance Function Estimation}

\begin{itemize}

\item Consider the partially linear model under high-dimensional regime
\begin{align*}
    Y &= \beta_1D+\eta(W)+\epsilon \\
    \beta_1
    &=
    \arg\min_b
    E\!\left[(\tilde Y-b\tilde D)^2\right] \\
    \tilde Y &= Y-E\!\left[Y\mid W\right] \\
    \tilde D &= D-E\!\left[D\mid W\right]
\end{align*}

    \item Here, $Y$ is product demand, $D$ is price, and $W$ is the BERT representation of the product description, obtained from the contextual embedding of the \texttt{[CLS]} token 
\end{itemize}
\end{frame}

\begin{frame}{Fine-Tuning BERT for Nuisance Function Estimation}

\begin{itemize}
\item To better approximate the nuisance functions
\[
g(W)=E\!\left[Y\mid W\right],
\qquad
m(W)=E\!\left[D\mid W\right]
\]
we adapt the pre-trained BERT model either by \textit{linear probing} or by \textit{fine-tuning}, which jointly updates both the BERT encoder and the prediction heads

\item We fine-tune BERT to learn product representations that better predict both price and demand, thereby more effectively adjusting for confounding information contained in product descriptions

\item Finally, we estimate price elasticity by applying Double Machine Learning to the residualized outcome $\tilde Y$ and treatment $\tilde D$ using cross-fitting

\end{itemize}
\end{frame}

\begin{frame}{Thank You}

\begin{itemize}
    \item Thank you! Please feel free to reach out with questions or to chat
    \textcolor{blue}{ \texttt{wenhao.jiang@duke.edu}}  \pause
    
    \begin{minipage}{0.7\textwidth}
        I am also looking for volunteers to help pre-test a research project that Siwei Cheng (NYU) and I are conducting on social interactions at conferences \\
        We do not collect your name, speech content, or mobility trajectory, only which badges you are close to
    \end{minipage} \hfill
    \begin{minipage}{0.22\textwidth}
        \begin{figure}
        \centering
        
\includegraphics[width=1\linewidth]{Feature Engineering/Figures/openrc28_qr.png}
    \end{figure}
    \end{minipage}
    
    \item If you are interested, please:
    \begin{itemize}
        \item Open the web app, add it to your home screen, collect a badge, input your badge number in your app, and wear it throughout the afternoon and Happy Hour
        \item At the end, return the badge to me or Cindy and confirm the return in the app
    \end{itemize}

\end{itemize}

\end{frame}

\begin{frame}{References}
    \footnotesize
    \bibliography{reference}
\end{frame}
\end{document}